% !TEX spellcheck = en_US
% !TEX spellcheck = LaTeX
\documentclass[a4paper,10pt,english]{article}
\usepackage{%
amsfonts,%
amsmath,%
amssymb,%
amsthm,%
babel,%
bbm,%
%biblatex,%
caption,%
centernot,%
color,%
enumerate,%
%enumitem,%
epsfig,%
epstopdf,%
etex,%
fancybox,%
framed,%
fullpage,%
%geometry,%
graphicx,%
hyperref,%
latexsym,%
mathptmx,%
mathtools,%
multicol,%
paralist,%
pgf,%
pgfplots,%
pgfplotstable,%
pgfpages,%
proof,%
psfrag,%
%subfigure,%
tcolorbox,%
tfrupee,%
tikz,%
times,%
%ulem,%
url,%
xcolor,%
mathpazo,%
}

\definecolor{shadecolor}{gray}{.95}%{rgb}{1,0,0}
\usepackage[margin=1in,top=0.75in]{geometry}
\usepackage[mathscr]{eucal}
\usepgflibrary{shapes}
\usepgfplotslibrary{fillbetween}
\usetikzlibrary{%
arrows,%
backgrounds,%
chains,%
decorations.pathmorphing,% /pgf/decoration/random steps | erste Graphik
decorations.text,% 
matrix,%
positioning,% wg. " of "
fit,%
patterns,%
petri,%
plotmarks,%
scopes,%
shadows,%
shapes.misc,% wg. rounded rectangle
shapes.arrows,%
shapes.callouts,%
shapes%
}

%\pgfplotsset{compat=newest} %<------ Here
\pgfplotsset{compat=1.11} %<------ Or use this one

\theoremstyle{plain}
\newtheorem{thm}{Theorem}[section]
\newtheorem{lem}[thm]{Lemma}
\newtheorem{prop}[thm]{Proposition}
\newtheorem{cor}[thm]{Corollary}
\newtheorem{clm}[thm]{Claim}

\theoremstyle{definition}
\newtheorem{defn}[thm]{Definition}
\newtheorem{conj}[thm]{Conjecture}
\newtheorem{exmp}[thm]{Example}
\newtheorem{axiom}[thm]{Axiom}
\newtheorem{assum}[thm]{Assumption}
\newtheorem{exerc}[thm]{Exercise}
\newtheorem*{defin}{Definition}

\theoremstyle{remark}
\newtheorem{rem}{Remark}
\newtheorem{note}{Note}

\newcommand{\iid}{\emph{i.i.d.}\ }
\newcommand{\Cov}{\operatorname{Cov}}
%\newcommand{\det}{\operatorname{det}}
%\newcommand{\Real}{\mathbb{R}}
\newcommand{\tr}{\operatorname{tr}}
\newcommand{\Var}{\operatorname{Var}}
\newcommand{\cov}{\operatorname{cov}}

%\renewcommand{\proof}[1]{\begin{proof}#1\end{proof}}
\newcommand{\EQ}[1]{\begin{equation*}#1\end{equation*}}
\newcommand{\EQN}[1]{\begin{equation}#1\end{equation}}
\newcommand{\eq}[1]{\begin{align*}#1\end{align*}}
\newcommand{\meq}[2]{\begin{xalignat*}{#1}#2\end{xalignat*}}
\newcommand{\norm}[1]{\left\lVert#1\right\rVert}
\newcommand{\inner}[1]{\left\langle #1\right\rangle}
\newcommand{\abs}[1]{\left\lvert#1\right\rvert}
\newcommand{\expect}[1]{\mathbb{E}\left[{#1}\right]}
\newcommand{\prob}[1]{\mathbb{P}\left[{#1}\right]}
\newcommand{\given}{\; \big\vert \;} 
\newcommand{\set}[1]{\left\{#1\right\}} 
\newcommand{\indicator}[1]{1_{\set{#1}}} 
\newcommand{\Ind}[1]{\mathbbm{1}_{#1}} 
\newcommand{\SetIn}[1]{\mathbbm{1}_{\set{#1}}} 
\newcommand{\red}[1]{\textcolor{red}{#1}} 
\newcommand{\floor}[1]{\left\lfloor#1\right\rfloor}
\newcommand{\ceil}[1]{\left\lceil#1\right\rceil}


\newcommand{\C}{\mathbb{C}}
\newcommand{\D}{\mathbb{D}}
\newcommand{\E}{\mathbb{E}}
\newcommand{\F}{\mathbb{F}}
\newcommand{\N}{\mathbb{N}}
\renewcommand{\P}{\mathbb{P}}
\newcommand{\Q}{\mathbb{Q}}
\newcommand{\R}{\mathbb{R}}
\newcommand{\Z}{\mathbb{Z}}

\newcommand{\bA}{\mathbf{A}}
\newcommand{\bI}{\mathbf{I}}
\newcommand{\bU}{\mathbf{1}}
\newcommand{\bx}{\mathbf{x}}
\newcommand{\bX}{\mathbf{X}}
\newcommand{\bZ}{\mathbf{Z}}


\newcommand{\cA}{\mathcal{A}}
\newcommand{\cB}{\mathcal{B}}
\newcommand{\cC}{\mathcal{C}}
\newcommand{\cF}{\mathcal{F}}
\newcommand{\cG}{\mathcal{G}}
\newcommand{\cM}{\mathcal{M}}
\newcommand{\cN}{\mathcal{N}}
\newcommand{\cP}{\mathcal{P}}
\newcommand{\cT}{\mathcal{T}}
\newcommand{\cX}{\mathcal{X}}
\newcommand{\cY}{\mathcal{Y}}

\newcommand{\sA}{\mathscr{A}}
\newcommand{\sB}{\mathscr{B}}
\newcommand{\sC}{\mathscr{C}}
\newcommand{\sE}{\mathscr{E}}
\newcommand{\sF}{\mathscr{F}}
\newcommand{\sG}{\mathscr{G}}
\newcommand{\sH}{\mathscr{H}}
\newcommand{\sL}{\mathscr{L}}
\newcommand{\sS}{\mathscr{S}}
\newcommand{\sT}{\mathscr{T}}
\newcommand{\sX}{\mathscr{X}}
\newcommand{\sY}{\mathscr{Y}}
\newcommand{\sZ}{\mathscr{Z}}

\newcommand{\tS}{\tilde{S}}
% Debug
\newcommand{\todo}[1]{\begin{color}{blue}{{\bf~[TODO:~#1]}}\end{color}}

% a few handy macros

\renewcommand{\le}{\leqslant}
\renewcommand{\ge}{\geqslant}
\newcommand\matlab{{\sc matlab}}
\newcommand{\goto}{\rightarrow}
\newcommand{\bigo}{{\mathcal O}}
%\newcommand{\half}{\frac{1}{2}}
%\newcommand\implies{\quad\Longrightarrow\quad}
\newcommand\reals{{{\rm l} \kern -.15em {\rm R} }}
\newcommand\complex{{\raisebox{.043ex}{\rule{0.07em}{1.56ex}} \hskip -.35em {\rm C}}}


% macros for matrices/vectors:

% matrix environment for vectors or matrices where elements are centered
\newenvironment{mat}{\left[\begin{array}{ccccccccccccccc}}{\end{array}\right]}
\newcommand\bcm{\begin{mat}}
\newcommand\ecm{\end{mat}}

% matrix environment for vectors or matrices where elements are right justifvied
\newenvironment{rmat}{\left[\begin{array}{rrrrrrrrrrrrr}}{\end{array}\right]}
\newcommand\brm{\begin{rmat}}
\newcommand\erm{\end{rmat}}

% for left brace and a set of choices
%\newenvironment{choices}{\left\{ \begin{array}{ll}}{\end{array}\right.}
\newcommand\when{&\text{if~}}
\newcommand\otherwise{&\text{otherwise}}
% sample usage:
%\delta_{ij} = \begin{choices} 1 \when i=j, \\ 0 \otherwise \end{choices}


% for labeling and referencing equations:
\newcommand{\eql}{\begin{equation}\label}
\newcommand{\eqn}[1]{(\ref{#1})}
% can then do
%\eql{eqnlabel}
%...
%\end{equation}
% and refer to it as equation \eqn{eqnlabel}.
% some useful macros for finite difference methods:
\newcommand\unp{U^{n+1}}
\newcommand\unm{U^{n-1}}
% for chemical reactions:
\newcommand{\react}[1]{\stackrel{K_{#1}}{\rightarrow}}
\newcommand{\reactb}[2]{\stackrel{K_{#1}}{~\stackrel{\rightleftharpoons}
 {\scriptstyle K_{#2}}}~}
\makeatletter
\def\th@plain{%
\thm@notefont{}% same as heading font
\itshape % body font
}
\def\th@definition{%
\thm@notefont{}% same as heading font
\normalfont % body font
}
\makeatother
\date{}


\title{Lecture-18: Stopping Times}
\author{}

\begin{document}
\maketitle

%\section{Moment Generating Function}
%\begin{defn} 
%For a random process $X:\Omega \to \sX^T$ defined on the probability space $(\Omega,\sF, P)$, 
%we can define the moment generating functional $\Psi_X: \sX^T \to \C$ as 
%\EQ{
%\Psi_X(f) = \E\exp(j \inner{f,X}),\quad\text{ for all } f \in \sX^T,
%}
%where $\inner{f,X}(\omega) = \sum_{t\in T}f_tX_t(\omega)$ for $T$ countable, 
%and $\inner{f,X}(\omega) = \int_{t\in T}f(t)X_t(\omega)dt$ for $T$ uncountable. 
%\end{defn}
%\begin{shaded}
%\end{shaded}

\section{Stopping times}
Let $(\Omega, \sF, P)$ be a probability space. %, and $T\subseteq \R$ an ordered index set considered as time. 
Consider a random process $X:\Omega\to\sX^T$ defined on this probability space with state space $\sX \subseteq \R$ and ordered index set $T\subseteq \R$ considered as time. 
\begin{defn}
A collection of event spaces denoted $\sG_{\bullet} \triangleq (\sG_t \subseteq \sF: t \in T)$ is called a filtration if $\sG_s \subseteq \sG_t$ for all $s \le t$. 
\end{defn}
\begin{rem} 
For the random process $X:\Omega\to\sX^T$, 
we can find the event space generated by all random variables until time $t$ as $\sF_t\triangleq \sigma(X_s, s \le t)$.  
The collection of event spaces $\sF_\bullet \triangleq (\sF_t: t \in T)$ is a filtration.
\end{rem}
\begin{defn}
The natural filtration associated with a random process $X:\Omega \to \sX^T$ is given by $\sF_\bullet \triangleq (\sF_t: t\in T)$ where $\sF_t \triangleq \sigma(X_s, s\le t)$.  
\end{defn}
\begin{rem}
For a random sequence $X: \Omega \to \sX^\N$, the natural filtration is a sequence $\sF_\bullet = (\sF_n: n \in \N)$ of event spaces $\sF_n \triangleq \sigma(X_1, \dots, X_n)$ for all $n \in \N$. 
\end{rem}
\begin{rem}
If the random sequence $X$ is independent, 
then the random sequence $(X_{n+j}: j \in \N)$ is independent of the event space $\sigma(X_1, \dots, X_n)$. 
\end{rem}
\begin{shaded}
\begin{exmp} 
For a random walk $S$ with step size sequence $X$, 
the natural filtration of the random walk is identical to that of the step size sequence. 
That is, $\sigma(S_1, \dots, S_n) = \sigma(X_1, \dots, X_n)$ for all $n \in \N$. 
This follows from the fact that for all $n \in \N$, 
we can can write $S_j = \sum_{i=1}^jX_i$ and $X_j = S_j - S_{j-1}$ for all $j \in [n]$. 
That is, there is a bijection between $(X_1, \dots, X_n)$ and $(S_1, \dots, S_n)$. 
\end{exmp}
\end{shaded}
\begin{defn}
\label{defn:StoppingTimeGen}
%For an ordered index set $T$, 
A random variable $\tau: \Omega \to T$ is called a \textbf{stopping time} with respect to  a filtration $\sF_\bullet$ if 
\begin{enumerate}[(a)]
\item the event $\tau^{-1}(-\infty, t] \in \sF_t$ for all $t \in T$, and 
\item the random variable $\tau$ is finite almost surely, i.e. $P\set{\tau < \infty} = 1$.  
\end{enumerate}
\end{defn}
\begin{rem}
%Consider the natural filtration $\sF_\bullet$ for a random process $X: \Omega \to \sX^T$,  defined by $\sF_t \triangleq \sigma(X_s, s \le t)$ for all $t\in T$. 
%We can consider the ordered index set $T$ as a time sequence. 
Intuitively, if we observe the process $X$ sequentially, 
then the event $\set{\tau \le t}$ can be completely determined by the observation $(X_s, s \le t)$ until time $t$. 
%denotes a stopping time for the process $X$. 
%Then, we have stopped after observing, $(X_s, s \le \tau)$, and before observing $(X_s, s > \tau)$.  
The intuition behind a stopping time is that its realization is determined by the past and present events but not by future events. 
That is, given the history of the process until time $t$, we can tell whether the stopping time is less than or  equal to $t$ or not. 
In particular, $\E[\SetIn{\tau \le t}\given \sF_t]=\SetIn{\tau \le t}$ is either one or zero. 
\end{rem}
\begin{defn}[First hitting time]
For a process $X:\Omega\to\sX^T$ and any Borel measurable set $A \in \sB(\sX)$, we define the \textbf{first hitting time} $\tau_X^A: \Omega \to T\cup\set{\infty}$ for the process $X$ to hit states in $A$, as 
$\tau_X^A \triangleq \inf\set{t \in T: X_t \in A}.$
\end{defn}
\begin{shaded}
%\begin{exmp}
%Consider a random bus trip, with each bus stop being a random output. 
%The random variable measuring ``time until bus crosses next stop after Majestic" is a stopping time as its value is determined by events before it happens. 
%On the other hand ``time until bus crosses the stop before Majestic'' would not be a stopping time  in the same context. 
%This is because we have to cross this stop, reach Majestic and then realize we have crossed that point. 
%%\item Consider $X_n \in \set{0,1}$ \textit{iid} Bernoulli$(1/2)$. Then $N = \min \set{n \in \N:\quad \sum_{i=1}^n X_i = 1}$ is a stopping time. 
%%For instance, $\Pr\set{N=2} = \Pr\set{X_1=0,X_2=1}$ and hence $N$ is a stopping time by definition. 
%\end{exmp}
\begin{exmp}
We observe that the event $\set{\tau_X^A \le t} = \set{X_s \in A \text{ for some }s \le t} \in\sF_t$ for all $t \in T$. 
It follows that, if $\tau_A$ is finite almost surely, then $\tau_A$ is a stopping time with respect to filtration $\sF_\bullet$.
\end{exmp}
\end{shaded}
\begin{prop} 
For a random sequence $X: \Omega \to \sX^\N$, 
an almost sure finite discrete random variable $\tau: \Omega \to \N\cup\set{\infty}$ is a \textbf{stopping time} with respect to this random sequence $X$ iff the event $\set{\tau = n} \in \sigma(X_1, \dots, X_n)$ for all $n \in \N$. 
%\begin{enumerate}[(i)]
%\item the event $\set{\tau = n} \in \sigma(X_1, \dots, X_n)$ for all $n \in \N$, and 
%\item the random variable is finite almost surely, i.e. $P\set{\tau < \infty}=1$. 
%\end{enumerate}
\end{prop}
\begin{proof}
From Definition~\ref{defn:StoppingTimeGen}, 
we have $\set{\tau = n} = \set{\tau \le n}\setminus\set{\tau \le n-1} \in \sigma(X_1, \dots, X_n)$. 
Conversely, from the theorem hypothesis, it follows that $\set{\tau \le n} = \cup_{m=1}^n\set{\tau = m} \in \sigma(X_1, \dots, X_n)$. 
\end{proof}
\begin{shaded}
\begin{exmp} 
Consider a random sequence $X: \Omega \to \sX^\N$, 
with the natural filtration $\sF_\bullet$,   
and a measurable set $A \in \sB(\sX)$. 
If the first hitting time $\tau_X^A:\Omega\to\N\cup\set{\infty}$ for the sequence $X$ to hit set $A$ is 
%defined by 
%\EQ{
%\tau_X^A \triangleq \inf\set{n \in \N: X_n \in A}.
%}  
%If $\tau_X^A$ is 
almost surely finite, then $\tau_X^A$ is a stopping time. 
This follows from the fact that $\set{\tau_X^A = n} = \cap_{k=1}^{n-1}\set{X_k \notin A}\cap\set{X_n \in A} \in \sF_n$ for each $n \in \N$. 
\end{exmp}
\end{shaded}
\begin{defn} 
Consider a random process $X:\Omega\to \sX^{\R_+}$ with discrete state space $\sX\subseteq\R$. 
For each state $y \in \sX$, we define $\tau_X^{\set{y},0} \triangleq 0$ and inductively define the \textbf{$k$th hitting time to a state $y$} after time $t = 0$, as 
\EQ{
\tau_X^{\set{y},k} \triangleq \inf\set{t > \tau_X^{\set{y}, k-1}: X_t = y},\quad k \in \N. 
} 
\end{defn}
\begin{rem} 
We observe that $\set{\tau_X^{\set{y}, k} \le t} \in \sF_t$ for all times $t \in \R_+$. 
Hence if $\tau_X^{\set{y}, k}$ is almost surely finite, then it is a stopping time for the process $X$. 
\end{rem}
\begin{defn} 
For a discrete valued random sequence $X:\Omega\to\sX^\N$, the number of visits to a state $y \in \sX$ in first $n$ time steps is defined as $N_y(n) \triangleq \sum_{k=1}^n\SetIn{X_k = y}$ for all $n \in \N$. 
\end{defn}
\begin{rem}
We observe that $N_y:\Omega\to\Z_+^{\Z_+}$ is a random walk with the Bernoulli step size sequence $(\SetIn{X_k = y}: k \in \N)$. 
Further, $\tau_X^{\set{y}, k} = \tau_{N_y}^{\set{k}} = \inf\set{n \in \N: N_y(n) = k}$. 
We also observe that $\{N_y(n) \le k\} = \{\tau_y^{(k+1)} > n\}$ and $\{N_y(n) = k\} = \{\tau_y^{k} \le n < \tau_y^{(k+1)}\}$. 
\end{rem}
\begin{rem}
We observe that the number of visits to state $y$ in first $n$ steps of $X$ is also given by 
\EQ{
N_y(n) 
= \sup\set{k \in \Z_+: \tau_X^{\set{y},k} \le n}
= \inf\set{k \in \N: \tau_X^{\set{y},k} > n} - 1
= \sum_{k\in\N}\SetIn{\tau_X^{\set{y},k} \le n}.
}
This implies that $N_y(n)+1$ is the first hitting time to set of states $\set{n+1, n+2, \dots }$ for the increasing random sequence $(\tau_X^{\set{y},k}:k \in \N)$. 
\end{rem}
\begin{lem}[Wald's Lemma] 
Consider a random walk $S: \Omega \to \R^{\Z_+}$ with \iid step-sizes $X: \Omega \to  \R^\N$ having finite $\E\abs{X_1}$. 
Let $\tau$ be a finite mean stopping time with respect to this random walk. 
Then,
\EQ{
\expect{S_\tau}= \expect{X_1}\expect{\tau}.
}
\end{lem}
\begin{proof} 
Recall that the event space generate by the random walk and the step-sizes are identical. 
From the independence of step sizes, it follows that $X_n$ is independent of $\sigma(X_0, X_1, \dots, X_{n-1})$. 
Since $\tau$ is a stopping time with respect to random walk $S$, 
we observe that $\set{\tau \ge n} = \set{\tau > n-1} \in \sigma(X_0, X_1, \dots, X_{n-1})$,  
and hence it follows that random variable $X_n$ and indicator $\SetIn{\tau \ge n}$ are independent and $\E[X_n\SetIn{\tau \ge n}] = \E X_1 \E\SetIn{\tau \ge n}$. 
Therefore,
\EQ{
\E[S_\tau]
= \E[\sum_{n=1}^\tau X_n]
= \E[\sum_{n \in \N} X_n\SetIn{\tau \ge n}]
= \sum_{n \in \N} \E X_1 \E\left[\SetIn{\tau \ge n}\right] 
= \E X_1\E\left[ \sum_{n \in \N}\SetIn{\tau \ge n}\right] 
= \E[X_1]\E[\tau].
}
We exchanged limit and expectation in the above step, which is not always allowed. 
We were able to do it %since the summand is positive and we apply monotone convergence theorem. 
by the application of dominated convergence theorem. 
%I'd like to point out here that in step (\ref{tricky}), you cannot always exchange infinite sums and expectations. 
%But here you can do so, because of the application of 
%Refer Ross/Wolff for more information. 
%Therefore, we can write
%	\begin{align*}
%	\sum_{n \in \N} \E\left[X_n 1_{\set{N \geq n}}\right] &= \sum_{n \in \N} \E\left[X_n\right]\E\left[ 1_{\set{N \geq n}}\right] \\
%	&= \E\left[X_1\right] \sum_{n \in \N} \Pr\set{N \geq n} \\
%	&= \E[X_1]\E[N].
%	\end{align*} 
%where the third equality follows from the fact that the expectation of a random variable being represented in terms of the \textit{ccdf} of the corresponding random variable.
\end{proof}
\begin{cor} 
Consider the stopping time $\tau_S^{\set{i}} \triangleq \min\set{n \in \N: S_n = i}$ for an integer random walk $S: \Omega \to \Z_+^{\Z_+}$ with \iid step size sequence $X: \Omega \to \Z^\N$.  
Then, the mean of stopping time $\E \tau_S^{\set{i}} = {i}/{\E X_1}$. 
\end{cor}
\begin{proof} 
This follows from the Wald's Lemma and the fact that $S_{\tau_i} = i$. 
\end{proof}


\subsection{Properties of stopping time} 

%Let the index set $T$ be separable, and 
\begin{lem}
Let $\tau_1,\tau_2$ be two stopping times with respect to filtration $\sF_\bullet$.  
Then the following hold true. 
\begin{enumerate}[i\_]
\item $\min \set{\tau_1,\tau_2} $ is a stopping time. 
\item If $T$ is separable, then $\tau_1+\tau_2$ is a stopping time.
\end{enumerate}
\end{lem}
\begin{proof}
Let $\tau_1,\tau_2$ be stopping times with respect to a filtration $\sF_{\bullet} = (\sF_t : t \in T)$.  
\begin{enumerate}[i\_]
\item Result follows since the event $\set{ \min \set{ \tau_1,\tau_2} > t} = \set{\tau_1 > t} \cap \set{\tau_2 > t} \in \sF_t$. 
%\EQ{
%\set{\tau_1\wedge\tau_2 > t} = \set{\tau_1 > t} \cap \set{\tau_2 > t} \in \sF_t.
%}
%The result follows since the events $\set{N_1 > n}$ and $\set{N_2 > n}$ depend solely on $\set{X_1, \dots, X_n}$. 
\item A topological space is called separable if it contains a countable dense set. 
Since $\R_+$ is separable and ordered, we assume $T = \R_+$ without any loss of generality. 
It suffices to show that the event $\set{\tau_1+\tau_2 \le  t} \in \sF_t$ for $T = \R_+$. 
To this end, we observe that 
\EQ{
\set{\tau_1 + \tau_2 \le t} = \bigcup_{s \in \Q_+:~ s \le t}\set{\tau_1 \le t - s, \tau_2 \le s} \in \sF_t.
}
\end{enumerate}
\end{proof}
\subsection{Strong independence property and applications}
\begin{thm}[Strong independence property] Let $X: \Omega\to\sX^{\R_+}$ be an independent random process with natural filtration $\sF_\bullet$, and $\tau:\Omega\to\R_+$ a stopping time with respect to $\sF_\bullet$, 
then $(X_{\tau+s}: s \in \R_+)$ is independent of history $(X_s: s \le \tau)$. 
\end{thm}
%\begin{proof}
%\end{proof}
\begin{defn} 
We can define the \textbf{$k$th return time to state $y$} for the random process $X:\Omega\to\sX^{\R_+}$ as the interval between two successive visits to state $y$, that is for all $k \in \N$
\EQ{
H_X^{\set{y}, k} 
\triangleq \tau_X^{\set{y},k} - \tau_X^{\set{y},k-1}
= \inf\set{s \in \R_+: X_{\tau_X^{\set{y},k-1}+s} = y}.
}
\end{defn}
\begin{rem}
We observe that $\tau_X^{\set{y}, k} = \sum_{j=1}^kH_X^{\set{y}, j}$ for all $k \in \N$. 
That is, the hitting times sequence $(\tau_X^{\set{y},k}: k \in \N)$ is a random walk with step-size sequence being return times $(H_X^{\set{y},j}: j \in \N)$. 
Therefore, if $H_X^{\set{y}, j}$ is almost sure finite for all $j \in [k]$, 
then the finite sum $\tau_X^{\set{y}, k}$ is almost sure finite for all $k\in \N$.  
\end{rem}
\begin{rem} 
From the bijection between hitting and return times, 
we observe that $\sigma(\tau_X^{\set{y}, k)}: k \in [n]) =  \sigma(H_X^{\set{y}, j}: j \in [n])$. 
Recall that $\tau_X^{N_y(n)+1} > n$ by definition. 
If passage times are independent and \iid from second passage time, 
then it follows from the Wald's Lemma that 
\EQ{
\E H_X^{\set{y},1} + \E[N_y(n)]\E H_X^{\set{y},2} \ge (n+1).
}
\end{rem}

\begin{shaded}
\begin{exmp}
We also observe that the $k$th hitting time to $\set{1}$ by a Bernoulli step size sequence $X:\Omega\to\set{0,1}^\N$ is the first hitting time to $\set{k}$ by random walk $S:\Omega\to\Z_+^{\Z_+}$. 
That is,  
\EQ{
\tau_X^{\set{1},k}
= \tau_S^{\set{k}} 
= \inf\set{n \in \N: S_n = k} 
%= \inf\set{n > \tau_S^{\set{k-1}}: S_n-(k-1) = 1}
= \tau_S^{\set{k-1}} + \inf\set{n \in \N: S_{\tau_S^{\set{k-1}}+n}-S_{\tau_S^{\set{k-1}}} = 1}. 
}
\end{exmp}
\end{shaded}
\begin{lem} 
For an \iid Bernoulli random sequence $X:\Omega\to\set{0,1}^\N$ with $\E X_1 \in (0,1)$, 
the $k$th hitting time to state~$1$ is a stopping time, and $\tau_X^{\set{1}, k} = \sum_{i=1}^kY_i$ where $Y:\Omega\to\N^\N$ is an \iid random sequence distributed identically to $\tau_X^{\set{1}}$ . 
\end{lem}
\begin{proof}
When the Bernoulli step size sequence $X$ is \iid with $\E X_1 = p \in (0,1)$, 
we get that $P\set{\tau_S^{\set{1}}= n} = (1-p)^{n-1}p$ for all $n \in \N$. 
It follows that 
\EQ{
P\set{\tau_S^{\set{1}}< \infty} 
= P\left(\cup_{n\in\N}\set{\tau_S^{\set{1}}=n}\right) 
= \sum_{n\in\N}P\set{\tau_S^{\set{1}}= n} 
= 1.
}
Hence, the random time $\tau_S^{\set{1}}$ is finite almost surely. 
We will show that $\tau_S^{\set{k}}$ is finite almost surely for all $k \in \N$ by induction. 
By induction hypothesis, $\tau_S^{\set{k-1}}$ is finite almost surely. 
Then $S_{\tau_S^{\set{k-1}}+n}-S_{\tau_S^{\set{k-1}}} = \sum_{j=1}^nX_{\tau_S^{\set{k-1}}+j}$ is the sum of $n$ \iid Bernoulli random variables independent of $\tau_S^{\set{k-1}}$ by strong independence property,  
and hence has distribution identical to $S_n$. 
Further,  
This implies that $\tau_S^{\set{k}} = \tau_S^{\set{k-1}} + \tau_S^{\set{1}}$, 
where $\tau_S^{\set{1}}$ has the identical distribution to $\tau_X^{\set{1}}$ and is independent of $\tau_S^{\set{k-1}}$. 
\end{proof}   
%Examples of stopping times.
%\begin{enumerate}
%\item For instance, \item Let $(N_n: n \in \N)$ be the number of successes for an \iid Bernoulli process $X$, then $T_k \triangleq \min\set{n \in \N: N_n = k}$ is a stopping time. 
%
%\item For any measurable set $A \in \sF$, the hitting time  $\min\set{n \in N: S_n \in A}$ of the set $A$ by random walk $S$ is a stopping time adapted to the natural filtration $\sF_{\bullet} = (\sF_n = \sigma(X_i, i \le n): n \in \N)$.
%%\item Let $S$ be a one-dimensional integer random walk with unit steps. 
%%% \textbf{Random walk:} 
%%% Consider $X_n$ \textit{iid} bivariate random variables with 
%%%\EQ{
%%%\Pr\set{X_n = 1} = \Pr\set{X_n = -1} = \frac{1}{2}. 
%%%}
%%Then $\tau_i \triangleq \min \set{n \in \N: S_n= i}$ is a stopping time.
%\end{enumerate}
\end{document}